\chapter{Spatial Orientation, Reorientation and Updating}
\label{sec:p19-spatial-orientation-and-updating}

Recognising an object after it has rotated is not the same task as
knowing where it lies after the observer has moved.
The identity benchmark in \cref{def:p19-rotation-benchmark} keeps
the comparison about objects; navigation can require a changing
relation between an observer, a remembered target and a surrounding
layout. This chapter develops that distinction through three bounded
experimental records: adult pointing after movement, adult transfer
between virtual room shapes, and children's search after disorientation
\cite{LiuXiao2018UpdatingSPU,AmbostaEtAl2013ReorientationSPU,
ShustermanEtAl2011NavigationDASHSPU}.
They do not measure a single interchangeable spatial ability.
The reporting standard is \cref{def:p19-evidence-record}.

\section{Frames, changes and questions}
\label{sec:p19-spu-frames}

\begin{hbdefinition}{def:p19-spu-tasks}{Three spatial questions}
\emph{Orientation} concerns an observer's heading or position relative
to a reference frame. \emph{Updating} concerns maintaining or revising
those relations as the observer moves, physically or in an instructed
imagined perspective. \emph{Reorientation} concerns recovering a usable
relation after orientation has been lost or made unreliable.
The tasks can interact, but success at one does not establish success
at all three. A pointing task can probe a remembered direction;
a corner-search task can probe recovery from ambiguity without
measuring an entire continuously updated trajectory.
\end{hbdefinition}

An egocentric relation uses an observer-centred reference: for example,
left of the current body heading, or a bearing relative to the head.
The relevant body or head reference must be specified.
An allocentric relation uses relations among environmental objects
or an environmental reference independent of the current observer
heading: the door is on the same wall as the window, or the target is
north of a stipulated map origin. A frame is not an assertion about
how long a memory lasts. An enduring representation can support an
egocentric response, and a briefly available environmental relation
can support an allocentric judgement. Likewise, a map, a verbal
instruction and a moving first-person display need not recruit the
same operations merely because all can express a direction.

\begin{hbexample}{def:p19-spu-turn-walk}{Turning is not translating}
Suppose an observer stands at a marked origin facing north.
One target is north of that origin and another is east of it.
After a right turn to face east without changing location, the
northern target is on the observer's left and the eastern target
is ahead. The target--target arrangement has not changed.
The observer's heading has.

Now suppose the observer walks east past the eastern target while
continuing to face east. That target is behind the observer; the
northern target's bearing also changes, but not by the same rule as
a pure turn at the origin. Recovering the original location and
recovering the original heading are separate demands.
Rotating a drawing of the two targets does not by itself solve the
translation problem. If the targets were only memorised, the answer
additionally depends on retaining their locations. Neither the
angle--time relation for object congruence nor the feature-binding
capacity result in \cref{sec:p19-mental-rotation} establishes how this
observer updates those relations.
\end{hbexample}

A response must also be separated from its interpretation.
Absolute angular error measures deviation in magnitude; signed error
retains the direction of a deviation. A common signed displacement
across targets is compatible with a wrong heading applied to an
otherwise usable configuration. Different target errors can instead
involve target memory, retrieval or response noise.
Neither pattern uniquely identifies a stored representation.
Response time includes comprehension, retrieval, transformation and
motor execution as well as any updating operation.

\section{Actual movement and imagined headings}
\label{sec:p19-spu-updating}

Liu and Xiao tested 120 university students, 60 men and 60 women,
aged 19--26, with normal or corrected-to-normal vision. Forty were assigned to each of
three environmental conditions; the report does not explicitly
establish random assignment. Four high-error participants were
excluded using an absolute-error threshold more than three standard
deviations above the mean. The analysed group sizes were 40, 38 and
38. The methods' wording about which retrieved group supplied two
exclusions is ambiguous; the analysis table supplies those final
denominators \cite[Methods and Table 1]{LiuXiao2018UpdatingSPU}.

Participants learned eight common objects arranged at equal intervals
on a one-metre radius, in a room partitioned by twelve wooden screens.
They practised the response task and repeatedly named and pointed to
the objects with their eyes closed during learning.
The \emph{immediate} group remained at the learning position.
The \emph{translated} group walked straight to a neighbouring novel
room, inspected it, and was blindfolded.
The \emph{disoriented} group followed a meandering walk until a
door-pointing error exceeded \(90^\circ\), before moving into the
neighbouring room. These conditions change access to the learned
environment and the history of movement, not just the amount of
mental rotation.

All participants were instructed to imagine the learning location
and heading. They physically turned left or right by \(90^\circ\);
the direction was counterbalanced. The test distinguished a
\emph{memory-aligned} imagined heading, matching the learning heading;
a \emph{sensorimotor-aligned} heading, matching the current body
heading; and a \emph{misaligned} heading matching neither.
Each heading required pointing to eight objects twice, giving 48
test trials per participant. Instructions were delivered through
earphones and pointing responses through a joystick.
Latency was measured from target-name onset to the response.
Participants were not continuously reporting their location along
the walking path \cite[Procedure]{LiuXiao2018UpdatingSPU}.

Both aligned headings had advantages over the misaligned heading.
The two aligned latencies did not differ detectably in the immediate
condition (\(p=.799\)) or translated condition (\(p=.119\)).
In the disoriented condition the memory-aligned response was faster
than the sensorimotor-aligned response
(\(F(1,37)=4.581\), \(p=.039\), reported partial
\(\eta^2=.11\)). Absolute error was smaller for the memory-aligned
than sensorimotor-aligned heading after translation
(\(F(1,37)=10.934\), \(p=.002\)) and disorientation
(\(F(1,37)=12.347\), \(p=.001\)), but not in the immediate condition
(\(F(1,39)=.083\), \(p=.775\)).
These non-significant differences are not equivalence tests.

The signed errors add information lost by a latency-only account.
The authors reversed the sign for one counterbalanced turning
direction to combine the conditions. Table~\ref{tab:p19-spu-bias}
reports their 95\% confidence intervals in degrees, not a new
calculation \cite[Results, pp.\ 5--6]{LiuXiao2018UpdatingSPU}.

\begin{table}[htbp]
\centering
\small
\caption{Signed pointing-error intervals in Liu and Xiao's experiment.
Memory alignment uses the learning heading; sensorimotor alignment
uses the current body heading.}
\label{tab:p19-spu-bias}
\begin{tabular}{lccc}
\hline
Condition & Memory aligned & Sensorimotor aligned & Misaligned\\
\hline
Immediate & \([-3.794,1.143]\) & \([4.961,9.297]\) & \([-10.502,-2.451]\)\\
Translated & \([-3.600,.951]\) & \([4.260,9.671]\) & \([-15.525,-4.953]\)\\
Disoriented & \([-1.589,2.841]\) & \([3.248,10.483]\) & \([-10.741,-4.149]\)\\
\hline
\end{tabular}
\end{table}

The memory-aligned intervals included zero; the sensorimotor-aligned
and misaligned intervals lay on opposite sides of zero in all three
conditions. Including zero does not establish exactly unbiased
pointing. The study's bar-chart captions describe intervals
corresponding to one standard error on either side of the mean;
those bars must not be substituted for the separately reported
95\% intervals in the table above.

The authors interpret the directional pattern as evidence for an
updating representation encoded independently of the long-term
memory representation. That interpretation uses a particular account
of bias during imagined-perspective transformation. It is not a
logical consequence of calling the account two systems.
Competition between accessible headings, inhibition of an irrelevant
heading, response remapping, and differences in retrieval availability
remain plausible contributors. The design constrains a single
latency interpretation by adding signed error, but it does not
identify independent neural modules. Its trained young-adult sample
and constrained room task do not establish accuracy during unrestricted
navigation, or effects of ageing or neurological injury.

\section{Recovering heading from environmental shape}
\label{sec:p19-spu-geometry}

A featureless rectangle has a practical ambiguity: two diagonally
opposite corners can share the same local arrangement of long and
short walls. A searcher can preserve this geometric relationship yet
select the rotationally equivalent wrong corner. Reorientation research
therefore distinguishes local corner geometry from properties of the
whole enclosure. Principal axes describe global axes of a shape;
medial axes form a branching internal skeleton sensitive to its
layout. Those are candidate descriptions of information used, not
names of proven brain mechanisms.

Ambosta, Reichert and Kelly recruited 132 first-year university
participants for course credit. The analysis retained 120 who
completed the task and passed training: 60 men and 60 women, mean
age 19.5. Twelve recruits failed training or did not complete the
task; the report does not give a complete breakdown of those
reasons. Participants were randomly assigned to learn a target
corner in one of two geometrically equivalent corner-pairs
\cite[Participants and Procedure]{AmbostaEtAl2013ReorientationSPU}.

This was desktop virtual navigation, using a mouse and keyboard,
not blindfolded physical locomotion. The training rectangle was
uniform grey, with a dark green floor and identical blue corner
response boxes. Its virtual dimensions corresponded to 20 by
5 metres. Training required a correct first choice on seven of
eight trials, with repeated blocks, feedback and retries, until
criterion or a 30-minute limit.
Testing comprised 48 trials: 24 reinforced rectangle-training
trials, 12 unreinforced rectangle controls, and 12 unreinforced
L-shaped-room transfer trials from six starting positions twice.
Unreinforced trials ended at the first choice.
The L-shaped room had five corner choices, allowing predictions
that coincide in the rectangle to diverge in the transfer task.

The source compared five specified principal-axis, local-geometry
and medial-axis prediction patterns. For an A/C-trained participant,
the principal-axis random account predicts choices across all five
corners; the nearest-corner account predicts A and D; and the
principal-axis sensorimotor account predicts U and D by combining an
axis-directed approach with a learned turn at a wall.
Local corner geometry predicts A and C, whereas the medial-axis
account predicts A, C and U. The transformed enclosure therefore
tests information that the training rectangle alone cannot separate
\cite[Fig.\ 2]{AmbostaEtAl2013ReorientationSPU}.
Participants favoured the
geometrically corresponding corners in the L-shaped room, but also
chose its extra corner, labelled U.
For the group trained to the A/C pair, the mean choice proportions
were A: .376, C: .342 and U: .183, each with reported standard error
.020; B and D attracted .036 and .062.
For the B/D-trained group, B: .310, D: .353 and U: .178
had standard errors .020, .030 and .020 respectively.
The three corners predicted by the medial-axis account together
attracted .901 (standard error .017) and .840 (standard error .024)
in the two groups, exceeding the uniform-choice benchmark .60
in the authors' comparisons (\(p<.001\)).
These are group proportions across repeated trials, not success
probabilities for every new navigator.

The extra-corner responses constrain a strict account using only
the two matching local corners. They do not uniquely establish
medial axes: a mixture of local and medial information and a
reward-sensitive excitation/inhibition account can both describe
aspects of the pattern.
One reported comparison needs an additional boundary.
Table 1 prints an Akaike information criterion value of
\(-193.929\) for the combined local-geometry/medial-axis model and
\(-117.839\) for the medial-axis excitation/inhibition model,
while the accompanying text claims that the criterion favours the
latter. With the usual smaller-is-preferred criterion these printed
values imply the opposite ranking.
This chapter retains the observed choices but does not treat that
inconsistent model-selection statement as settled
\cite[Fig.\ 5 and Table 1]{AmbostaEtAl2013ReorientationSPU}.

The analysis used a .01 significance threshold for several
comparisons; a printed \(p=.025\) is not a detected difference
under that threshold, nor proof of identical behaviour.
The reported lack of gender effects in this task does not establish
equal performance in all spatial tasks.
Transfer from a trained virtual rectangle to this L-shaped enclosure
cannot settle all environmental geometries, innate geometric
knowledge, or the effects of vestibular and locomotor cues available
in physical movement.

\section{Landmarks, task relevance and children's search}
\label{sec:p19-spu-landmarks}

A coloured wall can be a direct beacon when the target is beside it,
or an indirect landmark when the target is elsewhere but its relation
to that wall is remembered. Merely perceiving or naming the wall
does not show that this relation has been encoded and used.
Shusterman, Lee and Spelke compared different verbal cues in
four-year-old children's reorientation
\cite{ShustermanEtAl2011NavigationDASHSPU}.
The held version is the Harvard DASH author-form manuscript of
the article published in \emph{Cognition}, not the publisher's
final PDF. The procedural and numerical statements below refer
to that held version.

Eighty-eight children, mean age four years three months, range
three years eleven months to four years eight months, participated
across a control group and four experiments.
The rectangular enclosure measured 1.2 by 1.8 metres, with one
red short wall and otherwise neutral walls.
The child saw a sticker hidden in a corner, was blindfolded and
turned three or four times in the held manuscript's procedure,
and then searched. A door-pointing check led to more turning if
orientation remained adequate; about one trial per three children
was discarded and replaced.
Each child completed two four-trial blocks, with the same hiding
corner within each block and a corner on the other geometric diagonal
in the second, also changing proximity to the red wall.
Facing directions and wall relationships were
counterbalanced. The apparatus differed between Experiments 1--2
and Experiments 3--4, a limit on attributing every cross-experiment
difference solely to cue wording.

\begin{hbdefinition}{def:p19-spu-search-codes}{Geometry is not landmark discrimination}
The first searched corner is coded C when correct, R when it is
the diagonally opposite rotational equivalent, and N or F for
the remaining near or far errors.
The proportion C+R indexes choices consistent with the rectangular
geometry. The difference between the proportions C and R, C--R,
indexes discrimination between the geometrically equivalent corners.
It is not the proportion correct conditional on choosing C or R.
A high C+R can coexist with poor discrimination of the actual
target corner.
\end{hbdefinition}

\begin{table}[htbp]
\centering
\small
\caption{Cue contrasts in the held Shusterman--Lee--Spelke manuscript.
Each child supplied four trials per block; participant counts are
not trial counts. Numerical outcomes are reported first-choice
percentages from the selected results prose.}
\label{tab:p19-spu-cues}
\begin{tabular}{p{.23\textwidth}cp{.29\textwidth}p{.29\textwidth}}
\hline
Condition & \(n\) & Block sequence & Reported correct first searches\\
\hline
Control & 16 & No cue; no cue & 45\%; 45\%\\
Spatial cue, Exp.\ 1 & 16 & No cue; relational wall cue & 46\%; 77\%\\
Reversed order, Exp.\ 2 & 16 & Relational cue; no cue & 77\%; 70\%\\
Wall mention, Exp.\ 3 & 16 & No cue; task-irrelevant wall description & 42\%; 55\%\\
Functional cue, Exp.\ 4 & 24 & No cue; wall described as useful for finding sticker & 36\%; 72\%\\
\hline
\end{tabular}
\end{table}

In Experiment 1 the relational cue referred to the target's wall
relationship without using left/right terminology.
The block-by-group contrast for C--R against controls was
\(F(1,30)=4.94\), \(p<.05\).
Experiment 2 reversed cue order: reported performance persisted
when the cue was removed in the second short block.
Its two C--R block means did not differ detectably (\(p=.91\));
that is not an equivalence test or evidence of durable transfer
outside this setting.
In Experiment 3, describing the wall as pretty without making
it functionally relevant yielded no detected main or interaction
effect against the control group. An order/cue-target relationship
also affected the results. The finding is not that wall mention
can never help.

Experiment 4 gave the wall a task-relevant but nonspatial function:
it could help the child find the sticker.
The C--R block-by-group contrast, with age as a covariate, was
\(F(1,37)=11.80\), \(p=.001\); its within-group increase was
\(t(23)=7.09\), \(p<.001\).
For the subgroup cued to the red wall while the target was at a
neutral wall, the manuscript reports 69\% correct versus 38\%
in the comparable control subgroup, with a one-tailed
\(t(14)=2.24\), \(p<.03\).
The corresponding C--R subgroup comparison did not reach
significance. The indirect-landmark result therefore has a
smaller denominator and a different evidential boundary than
the overall functional-cue contrast.

The author-form manuscript also has numerical reporting
inconsistencies: some wall-split C--R values in its Table 1 do
not reconcile with overall C--R means in its results prose,
and one geometry paragraph gives conflicting percentages.
This chapter does not use those disputed means or silently
repair them. The selected contrast statistics and first-search
percentages above are explicitly reports from the results prose,
not independently recovered raw data.
The source supplies no confidence interval for those selected
contrasts. This limits numerical precision and leaves a duty to
check the final version and, if possible, the underlying data.

The experiments support a role for how language makes a landmark
relevant to the task, not just whether colour vocabulary is
available. They do not uniquely identify pragmatic interpretation,
attention, binding of geometry and colour, or memory encoding
as the causal internal step.
The first-search outcome depends jointly on remembering the target,
recovering orientation, interpreting the cue and selecting a
response. A small developmental laboratory sample and a
few immediate searches do not establish a general navigational
training prescription.

\begin{hbexample}{def:p19-spu-landmark-case}{An informative wall needs a remembered relation}
After disorientation in a rectangle, a child remembers that the
sticker was at a corner with the long wall on one side and the
short wall on the other. Two opposite corners satisfy that
description. Remembering which of those corners was next to
the red wall can break the ambiguity, provided the wall has
remained in the same place and the remembered relation is correct.
Saying that the red wall is useful does not itself state the
target corner: the child must still use an encoded relationship.
If the wall has been moved, an otherwise successful landmark
strategy can mislead. If the target relationship was never
remembered, seeing the wall after turning cannot recreate that
missing information. The geometry, cue-validity premise, memory
and online inference each make a separate contribution.
\end{hbexample}

\section{What the comparisons establish}
\label{sec:p19-spu-boundaries}

\begin{hbmethod}{def:p19-spu-diagnosis}{Keep the transformation and the evidence apart}
Specify whether location, heading or object orientation changes;
which cues remain available during movement and at response;
whether movement is bodily, mouse-controlled or only imagined;
and whether the target layout is visible or must be retrieved.
Record the actual response and its denominator: latency, signed
or absolute pointing error, first-corner search, or repeated-trial
choice proportion.
Then compare an updating account with retrieval, response
compatibility, landmark learning and strategy alternatives.
An imagined aligned heading can make retrieval easier without
proving an online updating mechanism; a motor-compatible response
can be faster without preserving more spatial information.
\end{hbmethod}

The three records constrain different spatial operations.
They show why environmental shape, landmark relevance and
movement history are not rival explanations at one fixed level:
they can supply different information and constrain different
operations. Behavioural dissociations license bounded functional
comparisons, not a deterministic brain-module claim.
A geometrically correct answer and an empirical explanation of
how a person reached it remain distinct; descriptive performance
differences are not by themselves normative instructions.

The exact allocated Oxford chapter is Hegarty and Stull's
\emph{Visuospatial Thinking}, Chapter 31, pages 606--630
\cite{HegartyStull2012VisuospatialSPU}.
The three original experiments presented here address particular
orientation, reorientation and updating tasks. Their designs and results
are not interchangeable with one another or with a broader synthesis
of visuospatial thinking.
